% Generated draft: requires mathematical and source review.
\documentclass[11pt]{article}
\usepackage[margin=25mm]{geometry}
\usepackage{fontspec}
\setmainfont{Latin Modern Roman}
\usepackage{amsmath,amssymb,mathtools,unicode-math}
\setmathfont{Latin Modern Math}
\usepackage{enumitem,needspace,xurl,hyperref}
\setlength{\parindent}{0pt}
\setlength{\parskip}{5pt}
\setlength{\emergencystretch}{3em}
\newcommand{\N}{\mathbb{N}}
\newcommand{\Z}{\mathbb{Z}}
\newcommand{\Q}{\mathbb{Q}}
\newcommand{\R}{\mathbb{R}}
\newcommand{\C}{\mathbb{C}}
\newcommand{\F}{\mathbb{F}}
\newcommand{\A}{\mathbb{A}}
\newcommand{\PP}{\mathbb{P}}
\newcommand{\E}{\mathbb{E}}
\newcommand{\eps}{\varepsilon}
\newcommand{\dd}{\,\mathrm{d}}
\begin{document}
\section*{UnsolvedMath ID: 800965}
\textbf{Source problem code:} catalogue entry
\subsection{Definitions and mathematical statement}
Let \(S^2\) be the two-dimensional sphere, and let \(g_{\mathrm{rd}}\) denote its round metric of constant Gaussian curvature \(1\), so that \(\operatorname{Area}_{g_{\mathrm{rd}}}(S^2)=4\pi\). Write \(\mathcal M\) for the space of smooth Riemannian metrics on \(S^2\); smoothness is the conventional interpretation used in this draft, since the supplied statement gives no explicit regularity class. For \(g\in\mathcal M\), a nonconstant closed geodesic is a nonconstant periodic map \(\gamma:\R/\Z\to S^2\) satisfying \(\nabla^g_{\dot\gamma}\dot\gamma=0\). Its length is
\[
L_g(\gamma)=\int_0^1\sqrt{g_{\gamma(t)}(\dot\gamma(t),\dot\gamma(t))}\,\dd t.
\]
The source asks whether every area-\(4\pi\) metric sufficiently close to \(g_{\mathrm{rd}}\) admits such a geodesic of length at most \(2\pi\), with equality allowed. More explicitly, once a topology \(\mathcal T\) on \(\mathcal M\) is specified, the proposed assertion is
\[
\exists\,U\in\mathcal T\quad g_{\mathrm{rd}}\in U\quad\text{and}\quad
\forall g\in U,\qquad
\operatorname{Area}_g(S^2)=4\pi
\ \Longrightarrow\ 
\exists\,\gamma\text{ a nonconstant closed }g\text{-geodesic with }L_g(\gamma)\leq2\pi.
\]
The supplied record does not specify \(\mathcal T\), so the meaning of sufficiently close remains unresolved; this draft does not select a topology.

The source also describes the question as asking whether \(\sigma\geq1/\pi\) near the round metric, but does not define \(\sigma\). If \(\ell(g)\) denotes the length of a shortest nonconstant closed \(g\)-geodesic and the intended convention is
\[
\sigma(g)=\frac{\operatorname{Area}_g(S^2)}{\ell(g)^2},
\]
then, under the area normalization, the stated bound becomes
\[
\ell(g)\leq2\pi\quad\Longleftrightarrow\quad\sigma(g)\geq\frac1\pi.
\]
This interpretation of \(\sigma\) is conditional and requires confirmation from the source context.
\subsection{Short English statement}
The question attributed in the supplied record to Babenko asks whether every Riemannian metric on \(S^2\) of area \(4\pi\), sufficiently close to the unit round metric, has a nonconstant closed geodesic of length at most \(2\pi\). The topology of closeness and the definition of the source's quantity \(\sigma\) are not supplied.
\subsection{Sources}
\begin{itemize}
\item Local UnsolvedMath registry, permanent ID \(800965\), exact source problem code \texttt{catalogue entry}; supplied title: \emph{Open Problems and Questions about Geodesics --- Question 46 (Babenko)}. The registry lists Burns and Matveev as proposers and gives the year 2019; the background records attribution/context to Babenko.
\item Burns and Matveev, \emph{Open Problems and Questions about Geodesics} (2019, as identified by the supplied record), Question 46. Supplied source URL: \url{https://arxiv.org/abs/1308.5417}. The record reports access on 2026-10-16 and extraction by \texttt{arxiv\_latex}; these are supplied provenance details, not a claim of independent source access or verification.
\item The supplied registry labels the problem \texttt{open}. This is a supplied classification, not an independently verified assessment of the literature or present mathematical status.
\end{itemize}
\end{document}
